% Auto-generated person page
\documentclass[10pt, a4paper]{book}
\usepackage[margin=1in]{geometry}
\usepackage{array}
\usepackage{xcolor}
\usepackage{ragged2e}
\usepackage{graphicx}
\newcommand{\birthsymbol}{\ensuremath{\star}}
\newcommand{\deathsymbol}{\ensuremath{\dagger}}

\providecommand{\recordchapter}{2}

\begin{document}
\thispagestyle{empty}

\recordchapter{Stefania Krysztofik}{Tomasz Krysztofik}

\begin{tabular}{@{}p{0.38\textwidth} p{0.58\textwidth}}
  \fbox{%
    \parbox[c][6cm][c]{\linewidth}{%
      \centering      
        {\includegraphics[width=\linewidth,height=5.8cm,keepaspectratio]{/Users/kadasie/PRIV_PRIV/GrampsBook/src/pages/pagetokr/assets/pictures/101d32301e2a3ae667f55985af06.png}}
    }%
  } &
  \begin{minipage}[t]{\linewidth}
    \vspace{0pt}
    \Huge \textbf{Tomasz Krysztofik}\\[0.3em]
    \small \textit{czeladnik garbarski, gospodarz, rolnik}\\[0.3em]

\begin{tabular}{@{}l l}
Ojciec: & {Kacper Krzysztofik} \\ 
Matka: & {Marianna Jarzemczyk} \\
\end{tabular}
  \end{minipage} \\ 
\end{tabular}

\noindent
\begin{tabular}{@{}l l l}
  \birthsymbol & \textbf{Data Narodzin:} & 1863.12.14, Płochocin, Rokitno \\ 
  \deathsymbol & \textbf{Data Śmierci: } & 1933 \\ 

\end{tabular}
\vspace{1.5em}

\begin{tabular}{@{}l l l l l l}
& \textbf{Data} & \textbf{Miejsce} & \textbf{Wiek} & \textbf{Opis} & \textbf{Osoba}\\
\hline
    \birthsymbol & 1863.12.14 & Płochocin, Rokitno & 0 dni & Narodziny & \\
\birthsymbol & 1864.12.19 & Białuty, Leszno & 1 lat & Narodziny (Partner) & Józefa Jędrzejewska\\
\birthsymbol & 1886 & Płochocin, Rokitno & 22 lat & Narodziny (Córka) & Ewa Krysztofik\\
\char"26AD & 1886.02.02 & Płochocin, Rokitno & 22 lat & Ślub & \\
\birthsymbol & 1888 & Płochocin, Rokitno & 24 lat & Narodziny (Syn) & Stanisław Krysztofik\\
\birthsymbol & 1891.08.29 & Warszawa & 27 lat & Narodziny (Córka) & Janina Marianna Puchalska\\
\deathsymbol & 1892 & Rokitno & 28 lat & Śmierć (Matka) & Marianna Jarzemczyk\\
\birthsymbol & 1893 & Nadarzyn & 29 lat & Narodziny (Córka) & Anna Krysztofik\\
\birthsymbol & 1897 & Nadarzyn & 33 lat & Narodziny (Córka) & Genowefa Krzysztofik\\
\birthsymbol & 1899.04.01 & Rusiec, Nadarzyn & 35 lat & Narodziny (Córka) & Stefania Krzystofik\\
\birthsymbol & 1901 & Nadarzyn & 37 lat & Narodziny (Syn) & Władysław Krysztofik\\
\birthsymbol & 1903 & Nadarzyn & 39 lat & Narodziny (Syn) & Józef Krysztofik\\
\birthsymbol & 1904 & Nadarzyn & 40 lat & Narodziny (Córka) & Marianna Krysztofik\\
\birthsymbol & 1907 & Nadarzyn & 43 lat & Narodziny (Syn) & Wacław Krysztofik\\
\char"26AD & 1908 & Rokitno & 44 lat & Ślub (Syn) & Stanisław Krysztofik\\
\char"26AD & 1909 & Rokitno & 45 lat & Ślub (Córka) & Ewa Krysztofik\\
\char"26AD & 1912 & Krosna, Rokitno & 48 lat & Ślub (Córka) & Janina Marianna Puchalska\\
\deathsymbol & 1918.06.04 & Gołaszew, Rokitno & 54 lat & Śmierć (Córka) & Stefania Krzystofik\\
\deathsymbol & 1924.06.06 &  & 60 lat & Śmierć (Córka) & Ewa Krysztofik\\
\deathsymbol & 1925 &  & 61 lat & Śmierć (Siostra) & Zofia Witkowska\\
\deathsymbol & 1933 &  & 69 lat & Śmierć & \\
\end{tabular}

\vspace{1.5em}
\newpage
        \begin{center}
        \includegraphics[width=\textwidth]{/Users/kadasie/PRIV_PRIV/GrampsBook/src/pages/pagetokr/assets/pictures/101c3a2ac2ec7bbde3f0dc5ba2f4.png}
        \newline
        {Narodziny Tomasza Krysztofika}
        \newline   \end{center} \newpage
        \begin{center}
        \includegraphics[width=\textwidth]{/Users/kadasie/PRIV_PRIV/GrampsBook/src/pages/pagetokr/assets/pictures/101c3a65ae2d1299f8697c309fd6.png}
        \newline
        {Akt Małżeństwa Tomasza i Józefy}
        \newline 
            \fbox{
            \parbox{\linewidth}{Nr 5
Krysztofik [tak zapisano tu nazwisko, w dalszej części Aktu jest Krzysztofik] Tomasz z Józefą Jędrzejewską
Płochocin

Działo się w Rokitnie, dnia 21 stycznia/2 lutego 1886 roku, o godzinie 2 po południu.

Wiadomo czynimy, że w obecności świadków, Józefa Zwierzchowskiego, 31 lat i Walentego Jarzembczyka, 27 lat mających, obu rolników z Płochocina, na dniu dzisiejszym zawarte zostało religijnie małżeństwo, między:

Tomaszem Krzysztofik, kawalerem, rolnikiem, urodzonym i zamieszkałym w Płochocinie, 22 lata mającym, synem żyjącego Kacpra i Marianny z domu Jarzembczyk, małżonków Krzysztofik z Płochocina

a

Józefą Jędrzejewską, panną, urodzoną w Białutach w parafii Leszno, zamieszkałą przy rodzicach w Płochocinie, 21 lat mającą, córką Stanisława i Bogumiły z Bieleckich, małżonków Jędrzejewskich z Płochocina.

Małżeństwo to poprzedziły 3 zapowiedzi, 5, 12 i 19/17, 24 i 31 stycznia bieżącego roku, wygłoszone w miejscowym kościele parafialnym.

Nowo zaślubieni oświadczyli, że nie zawarli między sobą umowy przedślubnej.

Religijna ceremonia zaślubin sprawowana była przez niżej podpisanego proboszcza parafii Rokitno.

Akt ten odczytano nowo zaślubionym i świadkom, niepiśmiennym, przez Nas tylko podpisany został.

Ks. Wł. Sędziakowski.}}
              \end{center} 


\end{document}
